% !TeX program = lualatex
% !TeX encoding = utf8
% !TeX root = ../Thermod.tex

%=========================================================
\Opensolutionfile{answer}[\currfilebase/\currfilebase-Answers]
\Writetofile{answer}{\protect\section*{\nameref*{\currfilebase}}}
\chapter{Принципи екстремумів і термодинамічної рівноваги}\label{\currfilebase}
\makeatletter
\def\input@path{{\currfilebase/}}
\makeatother
%=========================================================

% --------------------------------------------------------
\section{1}
% --------------------------------------------------------
Теорія термодинамічної рівноваги була розвинена Гіббсом на подобі теорії рівноваги механічних систем. Як відомо, в механічній системі рівноважному положенню відповідає екстремальне положення потенціальної енергії системи-мінімум, або еквівалентно-для виводу системи із положення рівноваги необхідно виконати над нею роботу. Розумно припустити, що роль потенціальної енергії в термодинаміці можуть виконувати термодинамічні потенціали $U$, $F$, $H$, $G$, що володіють, як ми встановили усіма властивостями потенціальної енергії. Щоб показати це будемо спиратись на другий закон термодинаміки, згідно з яким, кожна ізольована система еволюціонує до стану, відповідаючи максимуму ентропії. І насправді, внаслідок наявності незворотних процесів в ізольованій системі ($U = \text{соnst}$, $V = \text{соnst}$) виконується співвідношення:
\begin{equation}\label{7.1}
	dS\geq0
\end{equation}
Поки нарешті не буде досягнуто стану,при котрому подальша зміна $S$ неможлива. З цього випливає що у цьому стані $S = S_{max}$.

% --------------------------------------------------------
\section{2}
% --------------------------------------------------------

Зауважимо, що \eqref{7.1} випливає з рівнянь термодинаміки:
\begin{equation}\label{7.2}
	dS \geq \dfrac{1}{T}(dU + PdV)
\end{equation}
Якщо припустити, що $U$ і $V$ --- соnst. Аналогічно з \eqref{7.1} випливає, що у випадку постійних $V$ і $S$ умовою рівноваги термодинамічної системи буде мінімум внутрішньої енергії системи. $U_{min} \Rightarrow \delta U = 0,~~ (V,~S = \text{соnst})$

Насправді, з \eqref{7.1} випливає, що при $V$ і $S$ незмінних, неврівноважені процеси протікають так, що $dU<0$, до тих пір поки не буде досягнуто мінімуму. Зауважимо, що умова максимуму ентропії і мінімуму внутрішньої енергії було виведено з умови існування незворотних процесів, для котрих справджується \eqref{7.1}. З точки зору математики ці умови слід записати таким чином:
\begin{equation}\label{7.3}
	\delta S = 0; \quad \delta^2 S < 0,
\end{equation}
що гарантує опуклість функції $S=S(U,V)$ в точці $S=S_{max}$. Аналогічно для $U=U(S,V)$:
\begin{equation}\label{7.4}
	\delta U =0; \quad \delta^2 U < 0,
\end{equation}

що гарантує увігнутість функції $U(S,V)$ в точці $U=U_{min}$ і тим самим стан $U_{min}$ є стійким. Зміст запису $\delta S$ і $\delta U$ маємо розуміти так, що мова йде про віртуальні зміни при збереженні умов $(dU=dV=0)$ або $(dV=dS=0)$, тобто змінах зовнішніх і внутрішніх параметрів сумісних з накладеними умовами. Наприклад, збільшення $\delta T$ в одній частині і зменшення в іншій на $\delta T$.

% --------------------------------------------------------
\section{3}
% --------------------------------------------------------

Звернемося тепер до систем що знаходяться у термостаті ($T=\text{соnst}$) при постійному об'ємі. Характеристичною функцією для такої системи буде $F=F(T,V)$. Її диференціал буде:
\begin{equation}\label{7.5}
	dF = dU -TdS
\end{equation}
Слідуючи логіці наведених міркувань неврівноважені процеси можуть лише зменшувати $P$. Дійсно, в найзагальнішому випадку:
\begin{equation}\label{7.6}
	dF\leq dU-TdS; \quad dF\leq -SdT-PdV
\end{equation}
Звідси, рівноважне значення системи з $T$, $U=\text{const}$ відповідає мінімуму вільної енергії Гельмгольца:
\begin{equation}\label{7.7}
	\delta F = 0; \quad \delta^2 F >0
\end{equation}

В зв'язку з тим що спад $P$ дорівнює роботі здійснюваною системою в квазістатичному процесі, то з \eqref{7.2} випливає теорема про максимальну роботу: У процесах, температура початкового і кінцевого стану яких однакова, максимальну роботу, котру може виконати система над середовищем, можна отримати лише в оборотних процесах для яких:
\begin{equation}\label{7.8}
	A_{max} = \Delta F_{12}
\end{equation}

Ця властивість дала привід Гельмгольцу називати $P$ вільною енергією. Різниця її значень для двох станів 1 та 2 з однаковою температурою $T_0$, виражає ту частину зміни внутрішньої енергії, котра при ізотермічному оборотному процесі може піти на виконання роботи. Залишкова частина $U-F=TS$ має назву \enquote{зв'язкова енергія}. Найбільш змістовна частина теореми про максимальну роботу полягає у вимозі оборотності процесу. І насправді, розмірковуючи подібним чином отримуємо, що для адіабатно ізольованих систем $A_{max} = \Delta U$,
що, утім, являє собою наслідок першого закону термодинаміки, але тільки для оборотних процесів. Узагальнюючи теорему, робимо висновок(укладаємо(заключаємо)), що максимальна робота може бути отримана в оборотних процесах і для розширених систем, для котрих:
\begin{equation}\label{7.9a}\tag{7.9а}
	A_{max} = \Delta G, \text{ при } T = \text{const}
\end{equation}
\begin{equation}\label{7.9b}\tag{7.9б}
	A_{max} = \Delta H, \text{ при } S = \text{const}
\end{equation}

% --------------------------------------------------------
\section{4}
% --------------------------------------------------------


Умова \eqref{77.5а} може бути переписана як:
\begin{equation}\label{78.1}
\Delta G = A_{\max} \geq -\int_{1}^{2}VdP, \quad dG \leq -SdT + Vdp
\end{equation}

оскільки робота розширеної системи дорівнює $-\int_{1}^{2}VdP$. При постійному тиску з \eqref{78.1} випливає:
\begin{equation}\label{78.2}
\Delta G < 0
\end{equation}

для неврівноважених процесів (або навіть рівноважних, але що відбуваються при $T$ і $p = \text{const}$). Звідси, умова рівноваги системи, що знаходиться у термостаті при постійному тиску буде:
\begin{equation}\label{78.3}
\delta G = 0, \quad \delta^2 G > 0
\end{equation}

Розмірковуючи подібним чином для адіабатно ізольованої системи, що знаходиться при $p=\text{const}$, знайдемо умову рівноваги:
\begin{equation}\label{78.4}
\delta H = 0, \quad \delta^2 H > 0
\end{equation}

Умова постійності ентропії системи здається важкою для виконання. Проте такі системи існують. Розглянемо хімічну реакцію:
\begin{equation}
H_2 + Cl_2 = 2HCl
\end{equation}

Така кількість молів не змінюється. Ентропія газу дорівнює:
\begin{equation}
S = \nu\left(C_V \ln T + R \ln\frac{V}{\vartheta} + S_0\right),
\end{equation}
де $C_V$ і $S_0$ для двоатомних газів, що входять у реакцію однакові, а отже $S=\text{const}$ під час реакції. Якщо реакція проводиться у замкненому об'ємі при $V=\text{const}$, то в силу ідеальності газу, тиск буде залишатися постійним, якщо підтримувати постійну температуру. Оскільки реакція екзотермічна, то ми маємо відводити тепло, що виділяється. Звернемо увагу, що у даному випадку $\lambda_v = \lambda_p$.